\documentclass[10pt,a4paper]{amsart}
\usepackage[margin=19mm]{geometry}
\usepackage{amsmath,amssymb,mathtools}
\usepackage[scr=rsfs,cal=euler]{mathalfa}
\usepackage{xfrac}
\usepackage[unicode,hidelinks,bookmarks=false]{hyperref}
\newcommand{\ie}{i.e.\ }
\newcommand{\cf}{cf.\ }
\newcommand{\resp}{respectively\ }
\newcommand{\Subsection}{\subsection}
\providecommand{\nobreakdash}{\nobreak\hbox{-}}
\setlength{\emergencystretch}{2em}
\multlinegap0pt
\usepackage{bm}
\usepackage{bbm}
\usepackage{dsfont}
\usepackage{xspace}
\usepackage{cancel}
\usepackage{mathtools}
\usepackage{amsthm}
\makeatletter
\@ifundefined{Thm}{\newtheorem{Thm}{Thm}}{}
\@ifundefined{appendixlemma}{\newtheorem{appendixlemma}{Lemma}}{}
\@ifundefined{lemma}{\newtheorem{lemma}[Thm]{Lemma}}{}
\makeatother
\newcommand{\Tr}{\operatorname{Tr}}
\newcommand{\Var}{\operatorname{Var}}
\newcommand{\conjugate}{{\operatorname{^*}}}
\title[Spectrum of random centrosymmetric matrices; CLT and Circula]{Spectrum of random centrosymmetric matrices; CLT and Circular law}
\author{Indrajit Jana, Sunita Rani}
\date{Source: arXiv:2403.03513v3 (CC BY 4.0)}
\begin{document}
\maketitle

\noindent\emph{Source and licence.} Spectrum of random centrosymmetric matrices; CLT and Circular law. Version arXiv:2403.03513v3 (CC BY 4.0), distributed under the Creative Commons Attribution 4.0 International licence, as recorded in the arXiv licence metadata for this exact version and shown on the abstract page https://arxiv.org/abs/2403.03513v3. Full rights notice: licenses/arxiv-2403.03513-CC-BY-4.0.txt.\par\smallskip

\noindent\textit{Editorial scope.} Counted unit: the Lemma whose source label is \texttt{None} in the article (source0.tex lines 676--678, source label \texttt{None}), delivered together with the article's own complete original proof of it (source0.tex lines 680--785, one \texttt{proof} environment of 106 lines). No mathematics is written, repaired or re-derived: statement and proof are the article's own text line for line. Required context retained uncounted: eqn: event that T\_1, T\_2 are bounded (lines 373--375); eqn: definition of chi (lines 666--670); lemma:Norm bound (lines 1230--1236). Every definition, display and label of those lines is retained unchanged. Omitted from this package and not redistributed: 1--372, 376--665, 671--675, 679--679, 786--1229, 1237--1474. Nothing inside a retained excerpt is omitted or altered: every retained body is delivered exactly as the article's lines are written, so no omission receipt of any kind accompanies this package. Citations: the retained excerpts carry no citation command at all, so no bibliography entry of the article is transcribed and no reference is redistributed. Reference adaptation: none was needed and none was made; every reference inside the delivered unit resolves inside it, so no unresolved reference or citation is produced. Printed slips kept literal: no printed word, symbol, hypothesis or number of the counted statement or of its proof was altered. Layout and class: the article's own class is \texttt{amsart} with packages including amsmath, amsthm, amssymb, mathrsfs, ifthen, graphicx, enumerate, hyperref, float, tikz, enumitem, fullpage, xcolor, todonotes; the wrapper is the lane's pinned \texttt{amsart} editorial wrapper at 10pt on A4, which restates the theorem-like environments the retained text needs and adds the article's own macro definitions that the retained text needs (\texttt{\textbackslash Tr}, \texttt{\textbackslash Var}, \texttt{\textbackslash conjugate}), with extra wrapper packages (bm, bbm, dsfont, xspace, cancel, mathtools). The delivered numbering is the unit's own. Assets: no retained span contains a figure environment, a graphics-inclusion command, a PDF-page inclusion, a table or a TikZ picture; no image file is copied into this package, the package declares no asset dependency and no image file is redistributed with it. Licence: version arXiv:2403.03513v3 of this article is distributed under the Creative Commons Attribution 4.0 International licence, as recorded in the arXiv licence metadata for this exact version and shown on the abstract page https://arxiv.org/abs/2403.03513v3. A full rights notice is delivered at licenses/arxiv-2403.03513-CC-BY-4.0.txt. Characters: every retained excerpt is pure ASCII, so the delivered engine, XeLaTeX with its default Latin Modern text font, reports no missing character.
\begin{equation}
\Omega_{n} = \{\|T_1\| \leq \rho_1 \} \cap \{\|T_2\| \leq \rho_2 \} \label{eqn: event that T_1, T_2 are bounded},
\end{equation}

\begin{align}
    &\sum_{k=1}^{\lceil n/2 \rceil}\mathbb{E}_{k-1}[(W_{z,k}(T_1)+W_{z,k}(T_2))(W_{\bar{\eta},k}(T_1)+W_{\bar{\eta},k}(T_2))]\notag\\
    &= \frac{\partial^2}{\partial z \partial \bar{\eta}}\left[\frac{1}{n}\sum _{k=1}^{\lceil n/2 \rceil}\Gamma_{k}(z,\bar{\eta})\right]+O(n^{-1}) \notag\\
    &=: \frac{\partial^2}{\partial z \partial \bar{\eta}}\chi_{k}(z,\bar{\eta})+O(n^{-1})\label{eqn: definition of chi}.
\end{align}

\begin{appendixlemma}\label{lemma:Norm bound}
    Let $T$ be a  random matrix with $m \times m$ dimension having i.i.d. entries with mean zero and variance one. Moreover, assume that the entries satisfy the Poincar\'e inequality with constant $\alpha$. Then there exists $\rho\geq1$ such that 
    $$
    \mathbb{P}(\|T/\sqrt{m}\|>\rho/4+t) \leq K \exp\left(-\sqrt{\frac{\alpha m}{2}}t\right),\;\;\;\; \forall \;\; t>0,
    $$
    where $K>0$ is a universal constant.
\end{appendixlemma}

\begin{lemma}
    Let $\chi_{k}(z, \bar{\eta})$ be the same as defined in \eqref{eqn: definition of chi}, that is, $\chi_{k}(z, \bar{\eta}) = \frac{1}{n}\sum _{k=1}^{\lceil n/2 \rceil}\Gamma_{k}(z,\bar{\eta}).$ Then $\Var(\chi_{k}(z, \bar\eta))\to 0$ as $n\to\infty.$
\end{lemma}

\begin{proof}
Since $\{W_{z,k}(T_1)+W_{z,k}(T_2)\}_n$ is a martingale difference sequence, we have 
$$ 
\mathbb{E}[(W_{z,k}(T_1)+W_{z,k}(T_2))(W_{\bar{\eta},l}(T_1\conjugate)+W_{\bar{\eta},l}(T_2\conjugate))]=0 \text{ if } k \neq l. $$
This implies,
\begin{align}
    &\mathbb{E}\left\{\sum_{k=1}^{\lceil n/2 \rceil}\mathbb{E}_{k-1}[(W_{z,k}(T_1)+W_{z,k}(T_2))(W_{\bar{\eta},k}(T_1\conjugate)+W_{\bar{\eta},k}(T_2\conjugate))]\right\}\notag\\
    &= \mathbb{E}\left\{\left(\sum_{k=1}^{\lceil n/2 \rceil}(W_{z,k}(T_1)+W_{z,k}(T_2))\right)\left(\sum_{k=1}^{\lceil n/2 \rceil}(W_{\bar{\eta},k}(T_1\conjugate)+W_{\bar{\eta},k}(T_2\conjugate))\right)\right\}\notag\\
    &= \mathbb{E}\left\{\left(\Tr \hat{\mathcal{R}}_z^{\circ}(T_1)+\Tr \hat{\mathcal{R}}_z^{\circ}(T_2)\right)\left(\Tr \hat{\mathcal{R}}_\eta^{\circ}(T_1\conjugate)+\Tr \hat{\mathcal{R}}_\eta^{\circ}(T_2\conjugate)\right)\right\}\notag\\
    &= \mathbb{E}\left\{\Tr \hat{\mathcal{R}}_z^{\circ}(M) \Tr \hat{\mathcal{R}}_{\bar{\eta}}^{\circ}(M^*)\right\}.\notag
\end{align}
Let us rewrite $\Gamma_{k}(z,\bar{\eta})$ as
\begin{align}
    &\Gamma_{k}(z,\bar{\eta})\notag\\
    &=4e_k^t \mathbb{E}_k [\hat{\mathcal{R}}_{z}(T_1^k)]\mathbb{E}_k [\hat{\mathcal{R}}_{\bar{\eta}}(T_1^k\conjugate)]e_k+4e_k^t \mathbb{E}_k [\hat{\mathcal{R}}_{z}(T_2^k)]\mathbb{E}_k [\hat{\mathcal{R}}_{\bar{\eta}}(T_2^k\conjugate)]e_k\notag\\
    &:= \Gamma_{k}^1(z,\bar{\eta})+\Gamma_{k}^2(z,\bar{\eta}).\notag
\end{align}
Consequently, we may decompose
\begin{align}
    &\Var\left(\frac{1}{n}\Gamma_{k}(z,\bar{\eta})\right)\notag\\
    &=\Var\left[\frac{1}{n}\Gamma_{k}^1(z,\bar{\eta})+\frac{1}{n}\Gamma_{k}^2(z,\bar{\eta})\right] \notag\\
    &\leq 2\Var\left(\frac{1}{n}\Gamma_{k}^1(z,\bar{\eta})\right)+2\Var\left(\frac{1}{n}\Gamma_{k}^2(z,\bar{\eta})\right). \notag
\end{align}
Therefore, proving the lemma is equivalent to showing
\begin{align*}
    \Var\left(\frac{1}{n}\Gamma_{k}^1(z,\bar{\eta})\right)+\Var\left(\frac{1}{n}\Gamma_{k}^2(z,\bar{\eta})\right) \rightarrow 0.
\end{align*}
We show that $\Var\left(\frac{1}{n}\Gamma_{k}^1(z,\bar{\eta})\right) \rightarrow 0.$ We shall use the Poincar\'e inequality to get estimates on $\Var\left(\frac{1}{n}\Gamma_{k}^1(z,\bar{\eta})\right).$ This is shown in the equation \eqref{eqn: variance bound by poincare inequality}. However, since we are working on the event $\Omega_{n, k}$, where norms of the matrices are bounded, some terms may not be differentiable with respect to the matrix entries. To circumvent such a problem, we define a smoothing function as follows.

% Recall 
% $$\Gamma_{k}^1(z,\bar{\eta}) =4e_k^t \mathbb{E}_k [\mathcal{R}_{z}(T_1^k)]I\mathbb{E}_k [\mathcal{R}_{\bar{\eta}}(T_1^k\conjugate)]e_k. $$

Let $\epsilon_n= 4/n, \text{ and } \omega_{\epsilon_n,k}:\mathbb{R}^{\lceil n^2/4 \rceil} \rightarrow [0,1]$ be a smooth function such that 
\begin{align*}
&\omega_{\epsilon_n,k}|_{\{\|T_1^k\|\leq \rho_1\}} \equiv 1, \\
&\omega_{\epsilon_n,k}|_{\{\|T_1^k\|\geq \rho_1+\epsilon_n\}} \equiv 0,\\
&\left|\frac{\partial}{\partial x_{i}}\omega_{\epsilon_n,k}\right| \leq \frac{2}{\epsilon_n}, \;\;\;\; \forall \;\; 1\leq i \leq \lceil n^{2}/4\rceil.
\end{align*}
Let us define the smoothed version of $\mathcal{R}$ as $$\tilde{\mathcal{R}}_z(T_1^k)=\mathcal{R}_{z}(T_1^k)\omega_{\epsilon_n,k},$$
and
$$
\tilde{\Gamma}_{k}^1(z,\bar{\eta}) =4e_k^t \mathbb{E}_k [\tilde{\mathcal{R}}_z(T_1^k)]\mathbb{E}_k [\tilde{\mathcal{R}}_{\bar{\eta}}(T_1^k\conjugate)]e_k.
$$
Since the smoothed version and the original one match everywhere except the narrow annulus $\{\rho_1<\|T_1^k\|<\rho_1+\epsilon_n\}$, we can easily obtain
$$
|\Gamma_{k}^1(z,\bar{\eta})-\tilde{\Gamma}_{k}^1(z,\bar{\eta})| \leq\frac{K}{(\tau - \epsilon_n)^2}1_{\{\rho_1<\|T_1^k\|<\rho_1+\epsilon_n\}},
$$

where $K$ is some universal constant. Now, by Lemma \ref{lemma:Norm bound}, we have, 
\begin{align}
    \mathbb{E}[|\Gamma_{k}^1(z,\bar{\eta})-\tilde\Gamma_{k}^1(z,\bar{\eta})|^2] \leq \frac{K}{(\tau - \epsilon_n)^4}  \exp{\left(-\tau \sqrt{\frac{\alpha n}{2}}\frac{3\rho_1}{4}\right)}.\label{eqn: diff between Gamma and tilde Gamma}
\end{align}

Therefore the smoothing function $\omega$ did not alter the original quantities by significant amount. For the notational convenience, let us denote
\begin{align*}
    \Lambda &:= \mathbb{E}_k[\mathcal{R}_{z}(T_1^k)], & \Upsilon &:= \mathbb{E}_k[\mathcal{R}_{\eta}(T_2^k)^\conjugate],\\
    \tilde{\Lambda} &:= \mathbb{E}_k[\tilde{\mathcal{R}}_z(T_1^k)], &\tilde{\Upsilon} &:= \mathbb{E}_k[\tilde{\mathcal{R}}_\eta(T_2^k)^\conjugate].
\end{align*}

So, $\tilde{\Gamma}_{k}^1(z,\bar{\eta}) = \sum_{s=0}^{\lceil n/2 \rceil} \Tilde{\Lambda}_{ks}\Tilde{\Upsilon}_{sk}.$ Since $x_{ij}s$ satisfy Poincar\'e inequality, we have

\begin{align}
    \Var(\Tilde\Gamma_{k}^1(z,\bar{\eta})) \leq \frac{1}{n\alpha} \sum_{i,j=1}^{k-1}\left[\mathbb{E}\left|\frac{\partial \Tilde\Gamma_{k}^1(z,\bar{\eta})}{\partial t^{1}_{ij}}\right|^2 +\mathbb{E}\left|\frac{\partial \Tilde\Gamma_{k}^1(z,\bar{\eta})}{\partial \bar{t^{1}}_{ij}}\right|^2\right].\label{eqn: variance bound by poincare inequality}
\end{align}
The sum is divided by $n$ because the matrix $T_{1}$ is scaled by $1/\sqrt{n}$. Moreover, the sum stops at $k-1$ because $\Tilde\Gamma_{k}^1(z,\bar{\eta})$ is a constant for $i, j = k,k+1,\dots, \lceil n/2 \rceil$. On the other hand, we have
\begin{align*}
    \frac{\partial \mathcal{R}_{z}(T_1^k)_{ks}}{\partial t_{ij}^1} &= \mathcal{R}_{z}(T_1^k)_{ki} \mathcal{R}_{z}(T_1^k)_{js}\mathbf{1}_{\{j\neq k\}},& 
    \frac{\partial \overline{\mathcal{R}_{z}(T_1^k)_{ks}}}{\partial t_{ij}^1}&=0, 
\end{align*}

As a result,
\begin{align*}
    &\frac{\partial \tilde{\Lambda}_{k s}}{\partial t^{1}_{i j}}=\tilde{\Lambda}_{k i} \tilde{\Lambda}_{j s}\mathbf{1}_{\{j\neq k\}}+\Lambda_{k s} \frac{\partial \omega_{\epsilon_n, k}}{\partial t^{1}_{i j}} 1_{\left\{\rho_1<\|T_{1}^{(k)}\|<\rho_1+\epsilon_n\right\}}, \\
    &\frac{\partial \tilde{\Upsilon}_{s k}}{\partial t^{1}_{i j}}=\Upsilon_{s k} \frac{\partial \omega_{\epsilon_n, k}}{\partial t^{1}_{i j}} 1_{\left\{\rho_1<\|T_{1}^{(k)}\|<\rho_1+\epsilon_n\right\}}, \\
    &\frac{\partial\left(\tilde{\Gamma}_{k}^1(z, \bar{\eta})\right)}{\partial t^{1}_{i j}}=  \sum_{s=0} ^{\lceil n/2 \rceil} \tilde{\Lambda}_{k i} \tilde{\Lambda}_{j s} \tilde{\Upsilon}_{s k}+\sum_{s=0} ^{n/2}\left(\Lambda_{k s } \tilde\Upsilon_{s k}+\tilde{\Lambda}_{k s} \Upsilon_{s k}\right) \frac{\partial \omega_{\epsilon_n, k}}{\partial t^{1}_{i j}} 1_{\left\{\rho_1<\|T_{1}^{(k)}\|<\rho_1+\epsilon_n\right\}}.
\end{align*}

Let us denote $\tilde{y}_{j k}=\sum_{s=0} ^{\lceil n/2 \rceil} \tilde{\Lambda}_{j s} \tilde{\Upsilon}_{s k}$. Then using the facts that $\|\tilde{\Lambda}\|,\|\tilde{\Upsilon}\| \leq\left(\tau-\epsilon_n\right)^{-1}$, we have
\begin{align*}
    &\sum_{i,j=1}^{k-1} \left|\sum_{s=0} ^{\lceil n/2 \rceil} \tilde{\Lambda}_{k i} \tilde{\Lambda}_{j s} \tilde{\Upsilon}_{s k}\right|^2=\sum_{i,j=1}^{k-1} \left|\tilde{\Lambda}_{k i} \tilde{y}_{j k}\right|^2 \leq\left\|\tilde{\Lambda}_k\right\|_2^2\left\|\tilde{y}_k\right\|_2^2 \leq\left(\tau-\epsilon_n\right)^{-6}, \\
    &\sum_{i,j=1}^{k-1} \left|\sum_{s=0} ^{\lceil n/2 \rceil}\left(\Lambda_{ks} \tilde{\Upsilon}_{s k}+\tilde{\Lambda}_{ks} \Upsilon_{s k}\right) \frac{\partial \omega_{\epsilon_n, k}}{\partial a_{i j}} \mathbf{1}_{\left\{\left(\rho_1<\|T_1^k\|<\rho_1+\epsilon_n\right\}\right.}\right|^2 \\
  &\leq \frac{10}{\epsilon_n^2} \sum_{i,j=1}^{k-1} \left\|\tilde{\Lambda}_k\right\|_2^2\left\|\tilde{\Upsilon}_k\right\|_2^2 1_{\left\{\rho<\|T_{1}^{(k)}\|<\rho+\epsilon_n\right\}} \leq 10 n^4\left(\tau-\epsilon_n\right)^{-4} 1_{\left\{\rho_1<\|T_{1}^{(k)}\|<\rho_1+\epsilon_n\right\}} .
\end{align*}

Using the above estimates, \eqref{eqn: event that T_1, T_2 are bounded}, we have
\begin{align}
    &\Var\left(\frac{1}{n} \sum_{k=1}^{\lceil n/2 \rceil} \tilde{\Gamma}_{k}^1(z, \bar{\eta})\right) \notag\\
    &\leq \frac{1}{n} \sum_{k=1}^{\lceil n/2 \rceil} \Var\left(\tilde{\Gamma}_{k}^1(z, \bar{\eta})\right)\notag\\
    &\leq \frac{2}{\alpha n\left(\tau-\epsilon_n\right)^6}+\frac{K n^3}{\left(\tau-\epsilon_n\right)^4} \exp \left(-\sqrt{\frac{\alpha n}{2} }\frac{3 \rho_{1}}{4}\right).\notag
\end{align}

Therefore from \eqref{eqn: diff between Gamma and tilde Gamma}, we obtain
\begin{align}
    &  \Var\left(\frac{1}{n} \sum_{k=1}^{\lceil n/2 \rceil} \Gamma_{k}^1(z, \bar{\eta})\right) \notag\\
    &\leq \frac{1}{n} \sum_{k=1}^{\lceil n/2 \rceil} \mathbb{E}\left[\left|\Gamma_{k}^1(z, \bar{\eta})-\tilde{\Gamma}_{k}^1(z, \bar{\eta})\right|^2\right]+\Var\left(\frac{1}{n} \sum_{k=1}^{\lceil n/2 \rceil} \tilde{\Gamma}_{k}^1(z, \bar{\eta})\right) \notag\\
    &\leq \frac{2}{\alpha n\left(\tau-\epsilon_n\right)^6}+\frac{K n^3}{\left(\tau-\epsilon_n\right)^4} \exp \left(-\sqrt{\frac{\alpha n}{2} }\frac{3 \rho_{1}}{4}\right) \rightarrow 0 \text{ as } n\rightarrow \infty.\notag
\end{align}

Similarly, $ \Var\left(\frac{1}{n} \sum_{k=1}^{\lceil n/2 \rceil} \Gamma_{k}^2(z, \bar{\eta})\right) \rightarrow 0 \text{ as } n \rightarrow \infty.$ Thus,
\begin{align}
    &\Var(\chi_{k}(z, \bar{\eta}))\notag\\
    &\leq   2\Var\left(\frac{1}{n} \sum_{k=1}^{\lceil n/2 \rceil} \Gamma_{k}^1(z, \bar{\eta})\right)+2\Var\left(\frac{1}{n} \sum_{k=1}^{\lceil n/2 \rceil} \Gamma_{k}^2(z, \bar{\eta})\right) 
     \rightarrow 0, \text { as } n\rightarrow \infty .\notag
\end{align}

\end{proof}

\end{document}
